% Autogenerated translation of 2015-10-20-math.md by Texpad
% To stop this file being overwritten during the typeset process, please move or remove this header

\documentclass[12pt]{book}
\usepackage{graphicx}
\usepackage{fontspec}
\usepackage[utf8]{inputenc}
\usepackage[a4paper,left=.5in,right=.5in,top=.3in,bottom=0.3in]{geometry}
\setlength\parindent{0pt}
\setlength{\parskip}{\baselineskip}
\setmainfont{Helvetica Neue}
\usepackage{hyperref}
\pagestyle{plain}
\begin{document}

\hrule
layout: post
title: a post with math
date: 2015-10-20 11:12:00-0400

\section*{description: an example of a blog post with some math}

This theme supports rendering beautiful math in inline and display modes using \href{https://www.mathjax.org/}{MathJax 3} engine. You just need to surround your math expression with \texttt{\$\$}, like \texttt{\$\$ E = mc\textasciicircum{}2 \$\$}. If you leave it inside a paragraph, it will produce an inline expression, just like \$\$ E = mc\textasciicircum{}2 \$\$.

To use display mode, again surround your expression with \texttt{\$\$} and place it as a separate paragraph. Here is an example:

\$\$
sum\emph{\{k=1\}\textasciicircum{}infty $\vert$langle x, e}k rangle$\vert$\textasciicircum{}2 leq $\vert$x$\vert$\textasciicircum{}2
\$\$

You can also use \texttt{\textbackslash{}begin\{equation\}...\textbackslash{}end\{equation\}} instead of \texttt{\$\$} for display mode math.
MathJax will automatically number equations:

begin\{equation\}
label\{eq:caushy-shwarz\}
left( sum\emph{\{k=1\}\textasciicircum{}n a}k b\emph{k right)\textasciicircum{}2 leq left( sum}\{k=1\}\textasciicircum{}n a\emph{k\textasciicircum{}2 right) left( sum}\{k=1\}\textasciicircum{}n b\_k\textasciicircum{}2 right)
end\{equation\}

and by adding \texttt{\textbackslash{}label\{...\}} inside the equation environment, we can now refer to the equation using \texttt{\textbackslash{}eqref}.

Note that MathJax 3 is \href{https://docs.mathjax.org/en/latest/upgrading/whats-new-3.0.html}{a major re-write of MathJax} that brought a significant improvement to the loading and rendering speed, which is now \href{http://www.intmath.com/cg5/katex-mathjax-comparison.php}{on par with KaTeX}.

\end{document}
